\pagestyle{empty}

\section*{CD Használati útmutató}

A szakdolgozathoz mellékelt adathordozó tartalmazza a dolgozat PDF-változatát, a dolgozat LaTeX forráskódját, valamint a dolgozathoz kapcsolódó forráskódokat és dokumentációkat.

Az adathordozó főbb jegyzékei és fájljai a következők:

\begin{itemize}
    \item \path{dolgozat.pdf}: a szakdolgozat végleges, PDF formátumú változata.

    \item \path{thesis/}: a szakdolgozat LaTeX forrásfájljai. Ebben a jegyzékben találhatók a dolgozat fejezetei, a hivatkozásokhoz szükséges fájlok, valamint a PDF előállításához használt egyéb LaTeX állományok.

    \item \path{docs/}: a könyvtárhoz és a dolgozat témájához kapcsolódó dokumentációk. Ezek a fájlok a megvalósított megoldások, illetve a vizsgált technológiák megértését segítik.

    \item \path{code/}: a szakdolgozathoz kapcsolódó programkódok.

    \begin{itemize}
        \item \path{code/examples/}: példaprogamok, amelyek bemutatják a könyvtár használatát.

        \item \path{code/library/cpp}: a dolgozat keretében elkészített saját C++ könyvtár implementációja.
    \end{itemize}

    \item \path{README.md}: a projekt rövid leírása, a szakdolgozat témája, célkitűzései, valamint a repository felépítésének áttekintése.
\end{itemize}

A dolgozat témája az alkalmazásfejlesztői szintű erőforrás-izoláció vizsgálata Linux rendszeren, különös tekintettel a Landlock és a seccomp mechanizmusokra. A mellékelt forráskódok egy olyan saját C++ könyvtár megvalósítását tartalmazzák, amelynek célja ezen biztonsági mechanizmusok egyszerűbb és kényelmesebb használata alkalmazásfejlesztői szinten.

A programkód fordításához és futtatásához a megfelelő aljegyzékben található útmutató, illetve a projekt \path{README.md} fájlja ad további információt. A C++ implementáció CMake alapú fordítási rendszert használ, a tesztek pedig a GoogleTest keretrendszerrel készültek.

Az adathordozó tartalmának megtekintéséhez elegendő a fájlokat egy tetszőleges könyvtárba másolni. A dolgozat a \path{dolgozat.pdf} fájl megnyitásával olvasható, a forráskódok pedig a \path{code/} jegyzék alatt találhatók.
